\documentclass[11pt]{article}
\input{packages.tex}
\usepackage{ulem}
\DeclareMathOperator{\sign}{sign}
\DeclareMathOperator{\Imm}{Im}
\DeclareMathOperator{\DIF}{DIF}
\DeclareMathOperator{\rg}{rg}
\DeclareMathOperator{\dist}{dist}

% Настройка колонтитулов
\pagestyle{fancy}
\fancyhf{} % Очистить все поля
\fancyhead[L]{Поверь мне, Карлсон, не в пирогах счастье }
\fancyhead[R]{\href{ https://t.me/mallina_olga }{Автор}}
\fancyfoot[C]{\thepage}               % Справа внизу

\setlength{\headheight}{13.6pt}

\hypersetup{
    colorlinks=true,
    linkcolor=blue,
    citecolor=blue,
    filecolor=blue,
    urlcolor=blue,
}

\begin{document}

\begin{titlepage}
    \centering
    \vspace*{\fill}  % заполняет пространство сверху

    {\Huge Лекция 8. Теорема Гамильтона-Кэли. Линейные функции. Билинейные и квадратичные формы.
    \par}

    \vspace*{\fill}  % заполняет пространство снизу
\end{titlepage}

\setcounter{page}{1}

\section{Теорема Гамильтона-Кэли }

\begin{theorem}[Гамильтона-Кэли] Пусть $p(\lambda)$ --- характеристический многочлен матрицы $A$. Тогда $p(A)=0$.
    ($p(\lambda)= det (A-\lambda E))$ 
\end{theorem}

\begin{proof}
    Если $\lambda$ --- не является корнем характеристического многочлена, то $det(A-\lambda E) \neq 0$.А значит матрица $(A-\lambda E)$ имеет обратную.
    Найдем обратную матрицу:
    \[ (A-\lambda E)^{-1} = \frac{B(\lambda)}{det(A-\lambda E)} \]
    где $B(\lambda)$ --- матрица с элементами  $b_{ij}= (-1)^{i+j} M_{ji}$ ,где  $ M_{ji}$ --- минор порядка  $n-1$.\\
    Умножим это равенство на $(A-\lambda E) \det (A - \lambda E)$ справа и получим:
    \[ E \cdot det(A-\lambda E) = B(\lambda) (A-\lambda E) \]
    Матрицу $B(\lambda)$ можно представить в виде: 
    \[ B(\lambda) = B_0+ B_1\lambda + B_2 \lambda ^2 + ... + B_{n-1} \lambda^{n-1}  \]
    Пусть $p(\lambda) = a_0 + a_1\lambda + ... + a_n\lambda^n $. Тогда, вспомнив, что $det(A-\lambda E) = p(\lambda)$, получим:
    \[ a_0 E +a_1 \lambda E + ... + a_n \lambda ^n E =B(\lambda) (A-\lambda E) \]
    Получили равенство двух многочленов. При этом так как два многочлена равны, то равны их соответствующие коэффициенты. 
    Соберём коэффициенты перед степенями $\lambda$ справа и слева:
    \[ \text{перед } \lambda^0: a_0 E = B_0A \]
    \[ \text{перед } \lambda^1: a_1 E = B_1A-B_0 \]
    \[ \text{перед } \lambda^2: a_2 E = B_2A-B_1 \]
    \[ \vdots \]
    \[ \text{перед } \lambda^n: a_n E = -B_{n-1} \] 
    Умножим $a_1 E = B_1A-B_0$ на $A$, $a_2 E = B_2A-B_1$ --- на $A^2$ и т.д. 
    Сложим левую и правые части и получим: $p(A)=0$
\end{proof}

\section{Линейные функции}

\begin{definition}
    Пусть $L$ --- линейное пространство. $L \xrightarrow{f} \mathbb{R}$. Функцию $f(x)\in \mathbb{R}$ назовём линейной, если
    \[ \forall x,y\in L \ f(x+y) = f(x) + f(y) \]
    \[ \forall \alpha\in \mathbb{R}, \ \forall x\in L \ f(\alpha x) = \alpha f(x) \]
\end{definition}

\begin{remark}
    Из определения линейной функции мы понимаем, что линейная функция --- это частный случай линейного отображения, когда в качестве второго линейного пространства используется $\mathbb{R}$.
    Поэтому все теоремы, которые доказаны ранее про линейное отображение, работают и с линейными функциями. В частности рассмотрим матрицу линейного отображения:\\
    Пусть $e_1, ..., e_n$ --- базис, $x = (x_1,...,x_n)$. Тогда 
    \[ f(x) = (f(e_1), ... , f(e_n)) 
    \begin{pmatrix}
        x_1 \\
        \vdots \\
        x_n
    \end{pmatrix} \] 

\end{remark}

\begin{remark}
    (ещё одна теорема: как меняется матрица отображения при замене базиса). $A' = F^{-1}AS $
    Так как во множестве действительных чисел базис не меняется, то $F^{-1}= E$ и $A'= AS$.

\end{remark}


\begin{definition}
    Вспомним, что на множестве отображений можно ввести операции: сложение отображений и умножение на число. Пусть $L$ --- линейное пространство. Рассмотрим
    множество всех линейных функций, на котором зададим операции: сумма функций и умножение на число. В итоге получим линейное пространство, элементами которого являются функции. 
    Оно будет тоже являться n-мерным пространством. Его называют сопряжённым пространством и обозначают $L^*$.
    Выберем функции $f_1, ... , f_n$. Попробуем сделать так, чтобы 
    \[ f_i(e_j) = 
    \begin{cases}
        1, i=j\\
        0, i\neq j
    \end{cases} \]
    В итоге мы получим базис $f$ в $L^*$, который будем называть взаимным (биортогональным) по отношению к базису e в $L$.
\end{definition}

\section{Билинейные и квадратичные формы}

\begin{definition}
    Пусть $L$ --- линейное пространство. Билинейной функцией (билинейной формой) назовём такую функцию $b \colon L \times L \to \mathbb{R}$, которая упорядоченной паре $(x,y)$ ставит в соответствие действительное число $b(x, y)$,
    и при этом она должна быть линейна по каждому аргументу, то есть:
    \[ \forall x,y,z \in L \ b(x+y,z)= b(x,z)+b(y,z) \]
    \[ \forall x,y,z \in L \ b(z,x+y) = b(z,x) + b(z,y) \]
    \[ \forall x,y\in L, \forall \alpha\in \mathbb{R} \quad  b(\alpha x,y) = \alpha b(x,y) \]
    \[ \forall x,y\in L, \forall \alpha\in \mathbb{R} \quad  b(x,\alpha y) = \alpha b(x,y) \]

\end{definition}

\begin{definition}
    Запишем билинейную функцию в координатах:
    Пусть $x=(x_1,...,x_n)^T, y=(y_1,...,y_n)^T$, x и y --- столбцы $e_1,...,e_n$ --- базис. Тогда:
    \[ b(x,y) = b(x_1e_1+x_2e_2+...+x_ne_n, y_1e_1+y_2e_2+...+y_ne_n) =  \sum_{1\le i, j\le n} x_i y_j b(e_i,e_j) =\]
    \[ =(x_1,...,x_n)\begin{pmatrix}
        b(e_1,e_1) & \cdots & b(e_1,e_n) \\
        & \ddots & \vdots \\
        &        & b(e_n,e_n)
    \end{pmatrix}  
    \begin{pmatrix}
        y_1\\
        \vdots\\
        y_n
    \end{pmatrix}
    = x^T By \]
\end{definition}

\begin{definition}
    Матрицу, которая составлена из элементов, которые равны значению билинейной формы упорядоченной пары базисных векторов, будем называть матрицей билинейной формы и обозначать $B$.
\end{definition}

\begin{remark}
    (Как меняется матрица B при замене базиса). Пусть $e\xrightarrow{S} e'$ 
    \[ e'=eS\]
    \[ x= Sx', y=Sy' \]
    \[ b(x,y) = x^T By = (Sx')^T B Sy'= (x')^T S^T BSy' = (x')^T B'y' \]
    \[ \Rightarrow B' = S^T BS \]
\end{remark}

\begin{definition}
    Билинейная форма (билинейная функция) называется симметричной, если 
    \[ \forall x,y\in L \ b(x,y) = b(y,x) \]
\end{definition}

\begin{statement}
    $b(x,y)$ --- симметрична $\Leftrightarrow B$ --- симметрична. 
\end{statement}

\begin{proof}
   $ (\Rightarrow) \forall i,j \in 1,...n \hookrightarrow b(e_i,e_j)= b(e_j,e_i)$  в силу симметричности билинейной формы $\Rightarrow B$ --- симметрична.\\
   $ (\Leftarrow)$ Если матрица $B$ симметрична, то есть $B=B^T$, то 
   \[ b(x,y) = b^T(x,y) = (x^T By)^T = y^T B^T x = y^T Bx = b(y,x) \]
\end{proof}

\begin{theorem}
    (Следствия). $B'= S^T BS\quad  (S$ --- невырожденная)
    \[ 1) rg B'= rg B \]
    \[ 2) det B' = detB \cdot (detS)^2 \Rightarrow \text{либо определители матриц B и B' одновременно = 0, либо они одного знака} \]
\end{theorem}

\begin{example}
    (билинейной формы) $b(x,y) = x_1y_1 + 3x_1y_2- 5x_2y_1 + 11x_2y_2=$ 
    \[ = (x_1 x_2) \begin{pmatrix}
        1 &3  \\
        -5 & 11
    \end{pmatrix}
    \begin{pmatrix}
        y_1\\
        y_2
    \end{pmatrix}\]
\end{example}

\begin{definition}
    Пусть $L$ --- линейное пространство. Пусть в этом пространстве задана какая-то симметричная билинейная функция $b(x,y)$.
    Квадратичной формой будем называть 
    \[  k(x) = b(x,x) \]
\end{definition}

\begin{example}
    Пусть $ x_1y_1 + 6x_1y_2$ --- билинейная функция. Тогда 
    \[ x_1^2+6x_1x_2 = (x_1 x_2)
    \begin{pmatrix}
    1 & 3\\
    3 & 0
    \end{pmatrix} (x_1x_2) \text{--- квадратичная форма.} \]
    Также $x_1y_1+6x_2y_1, x_1y_1+2x_1y_2+4x_2y_1$ --- тоже билинейные функции, которыми может быть порождена данная квадратичная форма.
    Но из всех таких билинейных функций, которые порождают нашу квадратичную форму лишь одна уникальна --- симметричная : $x_1y_1 + 3x_1y_2 +3x_2y_1$.
\end{example}

\begin{proof}
    Покажем, что если есть квадратичная форма, то обязательно найдётся симметричная билинейная. Пусть $k(x)$ --- квадратичная форма.
    \[ k(x+y) = b(x+y, x+y) = b(x,x) + b(x,y) + b(y,x) + b(y,y) \]
    Но так как $b(x,y)$ --- симметричная, то $b(x,y) = b(y,x)$. Значит:
    \[ b(x,y) = \frac{k(x+y)-k(x)-k(y)}{2} \]
    $\Rightarrow b(x,y)$ по $k(x+y),k(x),k(y)$ определяется однозначно.   
\end{proof}

\begin{remark}
    $k(x) = x^TBx$
\end{remark}

\begin{definition}
    Будем говорить, что квадратичная форма $k(x)$ имеет диагональный вид, если
    \[ k(x) = \sum_{i=1}^n \alpha_i x_i^2 \]
\end{definition}

\begin{definition}
    Будем говорить, что квадратичная форма $k(x)$ имеет канонический вид, если
    \[ k(x) = \sum_{i=1}^n \alpha_i x_i^2,\ \alpha_i = 0, \pm 1 \]
\end{definition}


\begin{theorem}
    Для любой квадратичной формы $k(x)$ существует базис, в котором она имеет канонический вид.
\end{theorem}

\begin{proof}
    1 способ:\\
    Вспомним, что $B'= S^TBS$. Предположим, что $S$ --- это какая-то элементарная матрица. 
    Рассмотрим матрицу $B$:
    если $a_{11}\neq 0$, то вычтем первую строку из всех остальных с соответствующим коэффициентом, чтобы под $a_{11}$ стояли 0. 
    Эти операции не затронут первую строку, поэтому после того, как мы сделали операции со строчками, те же самые операции необходимо сделать со столбцами. Таким образом справа от $a_{11}$ будут стоять 0.
    \[ B=\left(
        \begin{array}{c|cccc}
            a_{11}\neq 0 & 0 & \cdots  &  0  \\ \hline
            0        \\
            \vdots   \\
            0      
        \end{array}
        \right)
        \]
    Затем переходим к следующей матрице. Действуем так до тех пор, пока не получим диагональный вид.
    Но сделать так можно не всегда, так как может встретиться особый случай, когда на диагонали появляется 0. 
    Если справа от него и под ним стоят 0, то просто идём дальше. Предположим, что 
    \[
    B=\left(
        \begin{array}{c|cccccc}
            a_{11}\ne 0 & 0 &  & \cdots &  & 0 & 0 \\ \hline
            0      & 0 &  0    & a\ne 0 &      \\
            \vdots & 0 &  &           \\
            0      & a\ne 0 &   \\
            \vdots   \\
            0         
        \end{array}
        \right) \]
    Далее необходимо посмотреть на элемент, который стоит на пересечении строки и столбца элементов $a\neq 0$. Если он не равен 0, то поменяем местами следующие строки:
    \begin{figure}[H]
        \centering
        \includegraphics[width=0.5\textwidth]{pictures/картиночка для понимания 1.png}
    \end{figure} 
    После того, как поменяли местами две строчки, необходимо поменять местами и соответствующие столбцы. В итоге получим, что на пересечении i-ой строки и j-ого столбца будет стоять ненулевой элемент, с которым можно проделать предыдущие операции. 
    Если же элемент на пересечении строки и столбца элементов $a\neq 0$ равен 0, то 
    \begin{figure}[H]
        \centering
        \includegraphics[width=0.5\textwidth]{pictures/еще одна картиночка для понимания.png}
    \end{figure}
    Итого, приведём квадратичную форму к диагональному виду. А из диагонального получим канонический:
    \[ \alpha_i x_i^2 = ( \sqrt {|\alpha_i|} x_i)^2 sign \ \alpha_i \]
    2 способ (метод Лагранжа выделения полных квадратов)\\
    Рассмотрим квадратичную форму $k(x)$ и будем считать, что $x_1^2$ присутствует в разложении формы: $k(x) = \alpha_1 (x_1^2 + \alpha_2' x_1 x_2+ \alpha_3' x_1 x_3 + ...   ) + \tilde{k}(x)$, где $\tilde{k}(x)$ не содержит $x_1$.
    \[ \alpha_2'x_1x_2 + \alpha_3' x_1 x_3 + ... \to + \left (\frac{\alpha_2' x_2 + ... + \alpha_n' x_n}{2}\right )^2 - \left (\frac{\alpha_2' x_2 + ... + \alpha_n' x_n}{2}\right )^2 \] 
    В итоге получим:
    \[ = \alpha_1(x_1+ ...) ^2 +  \tilde{\tilde{k}} \text{, где }\tilde{\tilde{k}} \text{ не содержит }x_1 \]
    Далее в оставшейся форме выделяем полный квадрат, например, с $x_2$ и так далее. 
    Выполним замену переменных:
    \[ x_1' = x_1 + ... \]
    \[ x_2' = 0\cdot x_1 + x_2 ... \]
    \[ \vdots\]
    \[ x_n' = x_n \]
    Матрица, соответствующая замене координат будет треугольной $\Rightarrow$ её определитель $\neq 0 \Rightarrow$ она является матрицей перехода. 
    Но может быть особый случай: квадраты закончились, а слагаемые без квадратов ещё есть. Допустим есть только $x_1x_2$, а их квадратов нет. Тогда делаем замену переменных:
    \[ x_1 = x_1'+ x_2'\]
    \[ x_2 = x_1' - x_2' \]
    \[ x_3= x_3' \]
    \[ \vdots \]
\end{proof}




\end{document}